\scanpage{285}
\mspage{13}

je renvoie à $\overline{\underline{\unsure{\mathrm{XIV}|\mathrm{XV}}}}$ \deleted{\unsure{8|3}}9.12 \inserted{(p. 133 bis).} Pour le nom ``axiome de connexité'', je peux offrir l'explication suivante 2.8. et 2.9.

\uline{Corollaire} 2.8. Si \deleted{\unsure{le} \unsure{est}} l'axiome de connexité $\operatorname{Loc}(5)$ est satisfait, alors les objets $W$-asphériques sont $0$-connexes (tautologique), et la réciproque est \unsure{vraie} si $W$ est fortement saturé, \deleted{ou si $\operatorname{Loc}(4b)$ est satisfait.}

\inserted{(ou aussi, si $W$ satisfait l'axiome des fibrations $\operatorname{Loc}\,\unsure{7}$ ci-dessous, puisque par $\overline{\underline{\mathrm{XII}}}$ th. 2 p. \unsure{21|31}, dans ce cas, $\operatorname{Loc}(5)$ équivaut à $W\ne\mathrm{Fl}(\unsure{M|\mathrm{Cat}})$ \unsure{\ldots})}

Car si l'axiome de connexité n'est pas satisfait, alors pour tt $X$ non vide \unsure{en|de} Cat, $X\unsure{{-}|\longrightarrow}e$ est iso, donc ($W$ étant fort\textsuperscript{mt} saturé) $X{-}e$ est dans $W$ i.e. $X$ est \deleted{dans} $W$-asphérique, donc il y a des objets \unsure{$W$|hot}{}-asphériques non $0$-connexes.

\deleted{Dans le cas où $\operatorname{Loc}(4b)$ est satisfait, et que tout objet $W$-asphé}

Notons aussi

\uline{Corollaire} 2.\unsure{9|8}. Supposons que $W$ satisfait à l'axiome $\operatorname{Loc}(4)$ des sommes discrètes (ou seulement à $\operatorname{Loc}\,4b)$. Pour que $W$ satisfasse l'axiome de connexité il suffit que pour tte flèche $f:X{-}Y$ avec $Y$ $0$-connexe, si $f$ est \deleted{\unsure{de}} $\in W$ alors $X$ est $0$-connexe. (C'est nécessaire bien sûr \unsure{sans|dans} \inserted{supposer} $\operatorname{Loc}(4b)$ ni rien \unsure{du|de} \deleted{\unsure{tout}}\unsure{tel|tout})

% archive sheet number: 284
